\documentclass[../../main.tex]{subfiles}
\begin{document}
% SOURCE p=186; printed=171; object=ch05-hilbert-spaces-title; region=lower
\section{Hilbert 空间}\label{section:ch05-hilbert-spaces}
% SOURCE p=186; printed=171; object=ch05-hilbert-introduction; region=lower
本节研究最重要的一类 Banach 空间，即 Hilbert 空间。在这种空间中，可以进行比一般 Banach 空间更精细的分析，而且许多结果直接推广了有限维 Euclid 几何中的熟悉事实。
% SOURCE p=186,187; printed=171,172; object=ch05-inner-product-definition; region=bottom-top
\begin{definition}[内积]\label{definition:ch05-inner-product}
复向量空间 $\mathscr H$ 上的\textbf{内积}，是从 $\mathscr H\times\mathscr H$ 到 $\mathbb C$ 的映射 $(x,y)\mapsto\langle x,y\rangle$，满足：
\begin{enumerate}
\item\label{item:ch05-inner-product-linearity} $\langle ax+by,z\rangle=a\langle x,z\rangle+b\langle y,z\rangle$，其中 $x,y,z\in\mathscr H$，$a,b\in\mathbb C$。
\item\label{item:ch05-inner-product-symmetry} $\langle y,x\rangle=\overline{\langle x,y\rangle}$。
\item\label{item:ch05-inner-product-positive} 当 $x\ne0$ 时，$\langle x,x\rangle\in(0,\infty)$。
\end{enumerate}
实向量空间上的内积类似定义，但取值和标量都为实数，第二条改为 $\langle y,x\rangle=\langle x,y\rangle$。
\end{definition}
% SOURCE p=186; printed=171; object=ch05-inner-product-second-linearity; region=bottom
由前两条，内积对第二个变量共轭线性：
\begin{align*}
\langle x,ay+bz\rangle=\overline a\langle x,y\rangle+\overline b\langle x,z\rangle.
\end{align*}
% SOURCE p=187; printed=172; object=ch05-prehilbert-definition; region=upper
\begin{definition}[预 Hilbert 空间与诱导范数]\label{definition:ch05-prehilbert}
带有内积的复向量空间称为\textbf{预 Hilbert 空间}。在这种空间中定义 $\|x\|=\sqrt{\langle x,x\rangle}$。
\end{definition}
% SUPPLEMENT: id=v14-ch05-inner-product-details; parent=ch05-inner-product-definition; kind=proof-detail+motivation+hypotheses; origin=independent-derivation
\begin{proof}
\textbf{基本恒等式补写（重构新增）。}第一条取零系数给 $\langle0,z\rangle=0$；第二条再给 $\langle z,0\rangle=0$。因此第三条与零向量情形合起来给非负性和正定性。共轭对称性与第一变量线性给
\begin{align*}
\langle x,ay+bz\rangle
=\overline{\langle ay+bz,x\rangle}
=\overline a\langle x,y\rangle+\overline b\langle x,z\rangle,
\end{align*}
特别地，$\langle ax,ax\rangle=|a|^2\langle x,x\rangle$。
\end{proof}
\begin{remark}
\textbf{动机与逐条件说明（重构新增）。}内积同时描述长度和夹角，并保留线性组合的代数信息。本书约定第一变量线性、第二变量共轭线性；另一种约定也常见，但不能在同一证明中混用。复数情形的共轭使平方长度实且非负；若两变量都复线性，则 $\langle\mathrm ix,\mathrm ix\rangle=-\langle x,x\rangle$，非零空间上不可能正定。正定性排除非零向量长度为零；删去它只能得到半内积和半范数。实数情形共轭不起作用。取平方根是为了让长度按标量的绝对值缩放；三角不等式尚须以下证明，不能由符号“范数”预先假定。预 Hilbert 空间暂不要求完备。
\end{remark}
% END SUPPLEMENT: id=v14-ch05-inner-product-details
% SOURCE p=187; printed=172; object=ch05-schwarz-theorem; number=5.19; region=upper
\begin{theorem}[Schwarz 不等式]\label{theorem:ch05-schwarz}
若 $x,y$ 属于一个预 Hilbert 空间，则 $|\langle x,y\rangle|\leqslant\|x\|\|y\|$；等号成立当且仅当 $x,y$ 线性相关。
\end{theorem}
% SOURCE p=187; printed=172; object=ch05-schwarz-proof; region=middle
\begin{proof}
当 $\langle x,y\rangle=0$ 时，结论显然。否则 $y\ne0$。令 $\alpha=\operatorname{sgn}\langle x,y\rangle$，$z=\alpha y$，则 $\langle x,z\rangle=\langle z,x\rangle=|\langle x,y\rangle|$ 且 $\|z\|=\|y\|$。对每个实数 $t$，
\begin{align*}
0\leqslant\|x-tz\|^2=\|x\|^2-2t|\langle x,y\rangle|+t^2\|y\|^2.
\end{align*}
右端二次多项式在 $t=\|y\|^{-2}|\langle x,y\rangle|$ 处达到最小值，故
\begin{align*}
0\leqslant\|x-tz\|^2=\|x\|^2-\|y\|^{-2}|\langle x,y\rangle|^2.
\end{align*}
不等式得证，而且等号成立恰在 $x=\alpha ty$ 时。
\end{proof}
% SUPPLEMENT: id=v14-ch05-schwarz-details; parent=ch05-schwarz-theorem; kind=proof-detail+motivation+hypotheses; origin=independent-derivation
\begin{proof}
\textbf{零值和等号条件（重构新增）。}内积非零时 $\alpha=\langle x,y\rangle/|\langle x,y\rangle|$，所以第二变量的共轭确实给 $\langle x,\alpha y\rangle=|\langle x,y\rangle|$。正定性使上式平方范数为零恰等价于 $x=\alpha ty$。若内积为零，等号成立恰当 $\|x\|\|y\|=0$，即至少一个向量为零，这时二者线性相关。反过来，若二者线性相关，含零向量时立即成立；否则可写 $x=cy$，于是 $|\langle x,y\rangle|=|c|\|y\|^2=\|x\|\|y\|$。
\end{proof}
\begin{remark}
\textbf{动机与条件说明（重构新增）。}比较 $x$ 与沿 $y$ 方向的所有实倍投影，把非负平方长度化为二次多项式，便得到控制混合项的基本界。相位 $\alpha$ 专门把复内积转成非负实数。正定性既保证除以 $\|y\|^2$ 合法，也负责等号对应向量相等；完备性完全未用。等号描述同一直线方向，而不是仅“正比例”，负数或复数比例都允许。
\end{remark}
% END SUPPLEMENT: id=v14-ch05-schwarz-details
% SOURCE p=187; printed=172; object=ch05-induced-norm-proposition; number=5.20; region=middle
\begin{proposition}\label{proposition:ch05-induced-norm}
在预 Hilbert 空间中，$x\mapsto\|x\|$ 是范数。
\end{proposition}
% SOURCE p=187; printed=172; object=ch05-induced-norm-proof; region=lower
\begin{proof}
$\|x\|=0$ 当且仅当 $x=0$，且 $\|ax\|=|a|\|x\|$，这两点显然。由 Schwarz 不等式，
\begin{align*}
\|x+y\|^2&=\|x\|^2+2\operatorname{Re}\langle x,y\rangle+\|y\|^2\\
&\leqslant\|x\|^2+2\|x\|\|y\|+\|y\|^2=(\|x\|+\|y\|)^2.
\end{align*}
这证明三角不等式。
\end{proof}
% SOURCE p=187; printed=172; object=ch05-hilbert-definition; region=lower
\begin{definition}[Hilbert 空间]\label{definition:ch05-hilbert-space}
关于上述范数完备的预 Hilbert 空间，称为\textbf{Hilbert 空间}。
\end{definition}
除非明确另作说明，“Hilbert 空间”总指复 Hilbert 空间；当然，也可以考虑实 Hilbert 空间。
% SUPPLEMENT: id=v14-ch05-hilbert-norm-details; parent=ch05-induced-norm-proposition; kind=proof-detail+motivation+hypotheses; origin=independent-derivation
\begin{remark}
\textbf{范数与完备性（重构新增）。}前面已证 $\langle ax,ax\rangle=|a|^2\langle x,x\rangle$，取非负平方根得到齐次性；正定性给零向量判别。原证最后两端非负，故可以取平方根，得到三角不等式。于是内积产生度量 $\|x-y\|$。Hilbert 定义中的完备性要求此度量的每个 Cauchy 列在原空间收敛，是以后构造最近点和无限正交展开的关键；它不能只换成代数上的无限维向量空间。例如有限支集序列空间 $c_{00}$ 配通常 $\ell^2$ 内积满足全部内积公理，却不完备：向量 $\sum_{j=1}^n2^{-j}e_j$ 为 Cauchy 列，其逐坐标极限支集无限，不在 $c_{00}$。实与复只是标量及共轭约定的区别，这些范数论证对两者都适用。
\end{remark}
% END SUPPLEMENT: id=v14-ch05-hilbert-norm-details
% SOURCE p=187,188; printed=172,173; object=ch05-l2-example; region=bottom-top
\begin{example}\label{example:ch05-l2-hilbert}
设 $(X,\mathscr M,\mu)$ 为测度空间。
\begin{definition*}[$L^2(\mu)$]\phantomsection\label{definition:ch05-l2-space}
以 $L^2(\mu)$ 表示满足 $\int|f|^2\,\mathrm d\mu<\infty$ 的可测复函数构成的空间，几乎处处相等的函数视为同一个元素。
\end{definition*}
不等式 $ab\leqslant\tfrac12(a^2+b^2)$ 给出 $|f\overline g|\leqslant\tfrac12(|f|^2+|g|^2)$，故 $f,g\in L^2(\mu)$ 时 $f\overline g\in L^1(\mu)$。由此，
\begin{align*}
\langle f,g\rangle=\int f\overline g\,\mathrm d\mu
\end{align*}
在 $L^2(\mu)$ 上定义一个内积。事实上，对于任意 $\mu$，$L^2(\mu)$ 都是 Hilbert 空间。我们将在\cref{theorem:ch06-lp-completeness}中证明完备性；目前先接受这一事实。
\begin{definition*}[$\ell^2(A)$]\phantomsection\label{definition:ch05-ell2-space}
若 $\mu$ 为 $(A,\mathscr P(A))$ 上的计数测度，其中 $A$ 为任意非空集合，相应的 $L^2$ 空间记作 $\ell^2(A)$。它由满足 $\sum_{\alpha\in A}|f(\alpha)|^2<\infty$ 的函数 $f:A\to\mathbb C$ 组成，这里的求和按\cref{section:ch00-extended-real-numbers}中的含义理解。
\end{definition*}
这一特殊情形的完备性也容易直接证明，见\cref{exercise:ch05-ell2-completeness}。
\end{example}
% SUPPLEMENT: id=v14-ch05-l2-completeness-details; parent=ch05-l2-example; kind=proof+motivation+hypotheses; origin=independent-derivation
\begin{proof}
\textbf{内积及完备性的完整补写（重构新增）。}不等式 $|af+bg|^2\leqslant2|a|^2|f|^2+2|b|^2|g|^2$ 保证 $L^2$ 对线性组合封闭。换取几乎处处相等的代表不改变乘积的积分，故内积良定义；线性、共轭对称性来自积分性质，而 $\int|f|^2=0$ 恰当 $f=0$ 几乎处处，给商空间上的正定性。前面的 Schwarz 不等式及范数命题现已适用，不需预先引用后面的 $L^p$ 理论。

设 $(f_n)$ 为 $L^2$ 中的 Cauchy 列。取严格递增 $n_k$，使 $\|f_{n_{k+1}}-f_{n_k}\|_2\leqslant2^{-k}$，并选可测有限值代表。令 $h_k=f_{n_{k+1}}-f_{n_k}$，$G_N=\sum_{k=1}^N|h_k|$。已证的三角不等式给 $\|G_N\|_2\leqslant\sum_{k=1}^N2^{-k}\leqslant1$。单调收敛定理给
\begin{align*}
\int\Bigl(\sum_{k=1}^{\infty}|h_k|\Bigr)^2\,\mathrm d\mu
=\lim_N\int G_N^2\,\mathrm d\mu\leqslant1.
\end{align*}
因此 $G=\sum_k|h_k|$ 有限几乎处处。令 $f=f_{n_1}+\sum_kh_k$ 于 $G<\infty$ 处，在可测零集 $\{G=\infty\}$ 上取0。它可测，且 $|f|\leqslant|f_{n_1}|+G$ 几乎处处，故 $f\in L^2$。对每个 $K$，同样对非负尾和应用单调收敛与有限项三角不等式，得
\begin{align*}
\|f-f_{n_K}\|_2
\leqslant\Bigl\|\sum_{k=K}^{\infty}|h_k|\Bigr\|_2
\leqslant\sum_{k=K}^{\infty}2^{-k}\longrightarrow0.
\end{align*}
任给 $\varepsilon>0$，先取 $N$ 使 $n,m\geqslant N$ 时 $\|f_n-f_m\|_2<\varepsilon/2$，再固定 $K$ 使 $n_K\geqslant N$ 且 $\|f_{n_K}-f\|_2<\varepsilon/2$，即有所有 $n\geqslant N$ 满足 $\|f_n-f\|_2<\varepsilon$。完备性得证。
\end{proof}
\begin{remark}
\textbf{动机与逐条件说明（重构新增）。}平方可积使“能量”有限，积分内积将有限维点积推广到函数。按几乎处处相等取商不可省，否则非零的零测集示性函数可能具有零范数。可测性保证积分有定义，平方可积通过原文的乘积界保证内积有限；不需要测度有限、$\sigma$ 有限或完备。补证只在一个可测零集上统一改值，因而不依赖测度完备。计数测度没有非空零测集，$\ell^2(A)$ 无须再识别不同函数；平方和定义为所有有限部分和的上确界，有限平方和保证非零坐标至多可数。$A$ 非空是原文约定；允许 $A=\varnothing$ 时得到零空间，同样完备。
\end{remark}
% END SUPPLEMENT: id=v14-ch05-l2-completeness-details
% SOURCE p=188; printed=173; object=ch05-hilbert-standing-assumption; region=upper
本节余下部分中，$\mathscr H$ 表示一个 Hilbert 空间。
% SOURCE p=188; printed=173; object=ch05-inner-product-continuity; number=5.21; region=upper
\begin{proposition}\label{proposition:ch05-inner-product-continuity}
若 $x_n\to x$ 且 $y_n\to y$，则 $\langle x_n,y_n\rangle\to\langle x,y\rangle$。
\end{proposition}
% SOURCE p=188; printed=173; object=ch05-inner-product-continuity-proof; region=upper
\begin{proof}
由 Schwarz 不等式，
\begin{align*}
|\langle x_n,y_n\rangle-\langle x,y\rangle|
&=|\langle x_n-x,y_n\rangle+\langle x,y_n-y\rangle|\\
&\leqslant\|x_n-x\|\|y_n\|+\|x\|\|y_n-y\|\longrightarrow0,
\end{align*}
因为 $\|y_n\|\to\|y\|$。
\end{proof}
% SUPPLEMENT: id=v14-ch05-inner-product-continuity-details; parent=ch05-inner-product-continuity; kind=proof-detail+motivation+hypotheses; origin=independent-derivation
\begin{remark}
\textbf{动机与条件说明（重构新增）。}两变量都变化时，先加减 $\langle x,y_n\rangle$，再分别控制两项。反三角不等式给 $|\|y_n\|-\|y\||\leqslant\|y_n-y\|$，故 $\|y_n\|$ 有界，原估计确实趋零。这证明内积的联合范数连续性；这里只需要预 Hilbert 空间，不用完备性。不能把范数收敛同时换成弱收敛：在 $\ell^2$ 中标准单位向量弱趋零，但 $\langle e_n,e_n\rangle=1$。
\end{remark}
% END SUPPLEMENT: id=v14-ch05-inner-product-continuity-details
% SOURCE p=188; printed=173; object=ch05-parallelogram-law; number=5.22; region=middle
\begin{theorem}[平行四边形恒等式]\label{theorem:ch05-parallelogram}
对任意 $x,y\in\mathscr H$，
\begin{align*}
\|x+y\|^2+\|x-y\|^2=2(\|x\|^2+\|y\|^2).
\end{align*}
\end{theorem}
这说明平行四边形两条对角线长度的平方和等于四条边长度的平方和。
% SOURCE p=188; printed=173; object=ch05-parallelogram-proof; region=middle
\begin{proof}
展开 $\|x\pm y\|^2=\|x\|^2\pm2\operatorname{Re}\langle x,y\rangle+\|y\|^2$，再把两式相加。
\end{proof}
% SUPPLEMENT: id=v14-ch05-parallelogram-details; parent=ch05-parallelogram-law; kind=motivation+hypotheses; origin=independent-derivation
\begin{remark}
\textbf{动机与条件说明（重构新增）。}相反符号的混合项抵消，留下纯长度关系；稍后它会把“两点都近乎达到最短距离”转成“两点彼此接近”。内积诱导范数是关键，完备性不参与此恒等式。任意范数不满足它，例如 $\mathbb R^2$ 的 $\ell^1$ 范数取 $x=(1,0),y=(0,1)$，左端为8，右端为4。
\end{remark}
% END SUPPLEMENT: id=v14-ch05-parallelogram-details
% SOURCE p=188; printed=173; object=ch05-orthogonality-definition; region=middle
\begin{definition}[正交与正交补]\label{definition:ch05-orthogonality}
若 $x,y\in\mathscr X$，当 $\langle x,y\rangle=0$ 时，称 $x$ 与 $y$ \textbf{正交}，记作 $x\perp y$。若 $E\subseteq\mathscr H$，定义
\begin{align*}
E^\perp=\{x\in\mathscr H:\langle x,y\rangle=0\text{ 对所有 }y\in E\}.
\end{align*}
\end{definition}
% SOURCE p=188; printed=173; object=ch05-orthogonal-complement-closed; region=middle
\begin{proposition}\label{proposition:ch05-orthogonal-complement-closed}
$E^\perp$ 是 $\mathscr H$ 的闭子空间。
\end{proposition}
% SOURCE p=188; printed=173; object=ch05-orthogonal-complement-closed-proof; region=middle
\begin{proof}
由\cref{proposition:ch05-inner-product-continuity}和内积的线性性立即得到。
\end{proof}
% SUPPLEMENT: id=v14-ch05-orthogonality-details; parent=ch05-orthogonality-definition; kind=proof-detail+source-clarification+motivation+hypotheses; origin=independent-derivation
\begin{remark}
\textbf{记号与动机（重构新增）。}原定义句中的 $\mathscr X$ 未另行指定；这里应理解为本节固定的 $\mathscr H$。正交把垂直关系推广为内积为零，共轭对称性保证 $x\perp y$ 当且仅当 $y\perp x$。正交补要求同时消去 $E$ 中每个方向；$E$ 不必是子空间或闭集，空集时 $E^\perp=\mathscr H$。以下闭性只用内积连续，不需要空间完备。
\end{remark}
\begin{proof}
\textbf{闭子空间结论补写（重构新增）。}对每个 $y\in E$，映射 $x\mapsto\langle x,y\rangle$ 线性连续，因此其核为闭子空间。$E^\perp$ 是这些核之交：交仍含0、对线性组合封闭，并且任意闭集之交闭。这也涵盖 $E$ 为空的情形。
\end{proof}
% END SUPPLEMENT: id=v14-ch05-orthogonality-details
% SOURCE p=188; printed=173; object=ch05-pythagorean-theorem; number=5.23; region=lower
\begin{theorem}[Pythagoras 定理]\label{theorem:ch05-pythagorean}
若 $x_1,\ldots,x_n\in\mathscr H$ 且 $j\ne k$ 时 $x_j\perp x_k$，则
\begin{align*}
\Bigl\|\sum_{j=1}^nx_j\Bigr\|^2=\sum_{j=1}^n\|x_j\|^2.
\end{align*}
\end{theorem}
% SOURCE p=188; printed=173; object=ch05-pythagorean-proof; region=lower
\begin{proof}
展开平方，
\begin{align*}
\Bigl\langle\sum_{j=1}^nx_j,\sum_{k=1}^nx_k\Bigr\rangle
=\sum_{j,k=1}^n\langle x_j,x_k\rangle=\sum_{j=1}^n\|x_j\|^2.
\end{align*}
\end{proof}
% SUPPLEMENT: id=v14-ch05-pythagorean-details; parent=ch05-pythagorean-theorem; kind=motivation+hypotheses; origin=independent-derivation
\begin{remark}
\textbf{动机与逐条件说明（重构新增）。}两两正交使每个不同指标的交叉项都为零，从而总能量等于各方向能量之和。这里只是有限和，展开不涉及极限和完备性；无限和须另证收敛。两两正交是充分条件，不是这一个长度等式的必要条件：复空间中 $x_1=1,x_2=\mathrm i$ 的实混合项为零，等式成立但内积不为零。若毫无正交控制，则 $x_1=x_2\ne0$ 就使等式失败。允许某些向量为零；若容许空族，空和也满足等式。
\end{remark}
% END SUPPLEMENT: id=v14-ch05-pythagorean-details
% SOURCE p=188; printed=173; object=ch05-orthogonal-decomposition; number=5.24; region=bottom
\begin{theorem}\label{theorem:ch05-orthogonal-decomposition}
若 $\mathscr M$ 为 $\mathscr H$ 的闭子空间，则 $\mathscr H=\mathscr M\oplus\mathscr M^\perp$；即每个 $x\in\mathscr H$ 有唯一分解 $x=y+z$，其中 $y\in\mathscr M$，$z\in\mathscr M^\perp$。此外，$y,z$ 分别是 $\mathscr M,\mathscr M^\perp$ 中距离 $x$ 最近的唯一点。
\end{theorem}
% SOURCE p=189; printed=174; object=ch05-orthogonal-decomposition-proof; region=upper
\begin{proof}
令 $\delta=\inf\{\|x-y\|:y\in\mathscr M\}$，并选 $\{y_n\}\subseteq\mathscr M$ 使 $\|x-y_n\|\to\delta$。由平行四边形恒等式，
\begin{align*}
2(\|y_n-x\|^2+\|y_m-x\|^2)
=\|y_n-y_m\|^2+\|y_n+y_m-2x\|^2.
\end{align*}
由于 $\tfrac12(y_n+y_m)\in\mathscr M$，最后一项至少为 $4\delta^2$，所以
\begin{align*}
\|y_n-y_m\|^2\leqslant2(\|y_n-x\|^2+\|y_m-x\|^2)-4\delta^2\longrightarrow0.
\end{align*}
因此 $(y_n)$ 为 Cauchy 列；由完备性和闭性，它趋于某个 $y\in\mathscr M$，且 $\|x-y\|=\delta$。

令 $z=x-y$。给定 $u\in\mathscr M$，将它乘适当非零标量后，可假定 $\langle z,u\rangle$ 为实数。函数
\begin{align*}
f(t)=\|z+tu\|^2=\|z\|^2+2t\langle z,u\rangle+t^2\|u\|^2\qquad(t\in\mathbb R)
\end{align*}
在 $t=0$ 处取最小值，故 $0=f'(0)=2\langle z,u\rangle$。因此 $z\in\mathscr M^\perp$。

若 $z'\in\mathscr M^\perp$，由 Pythagoras 定理，
\begin{align*}
\|x-z'\|^2=\|x-z\|^2+\|z-z'\|^2\geqslant\|x-z\|^2,
\end{align*}
等号恰在 $z'=z$ 时成立。因此 $z$ 是相应的唯一最近点，同样的论证给出 $y$ 的唯一性。分解的唯一性也由 $\mathscr M\cap\mathscr M^\perp=\{0\}$ 得到。
\end{proof}
% SUPPLEMENT: id=v14-ch05-orthogonal-decomposition-details; parent=ch05-orthogonal-decomposition; kind=proof-detail+motivation+hypotheses; origin=independent-derivation
\begin{proof}
\textbf{极小化与相位细节（重构新增）。}$0\in\mathscr M$，故 $0\leqslant\delta\leqslant\|x\|$，按下确界定义可取 $\|x-y_n\|<\delta+1/n$。范数连续给极限点达到距离。对 $u\in\mathscr M$，若 $\langle z,u\rangle\ne0$，取 $\alpha=\langle z,u\rangle/|\langle z,u\rangle|$；则 $\alpha u\in\mathscr M$ 且 $\langle z,\alpha u\rangle=|\langle z,u\rangle|$ 为实数。原证应用于 $\alpha u$ 得此值为0，矛盾。因此原来的内积必为零。其最小性来自 $z+tu=x-(y-tu)$ 且 $y-tu\in\mathscr M$。

对任意 $y'\in\mathscr M$，$z\perp y-y'$，故 $\|x-y'\|^2=\|z\|^2+\|y-y'\|^2$，等号达到最短距离恰当 $y'=y$。若还有另一分解，两分解之差同时属于 $\mathscr M$ 和 $\mathscr M^\perp$；这样的向量 $v$ 满足 $\|v\|^2=0$，故为零。
\end{proof}
\begin{remark}
\textbf{动机与条件说明（重构新增）。}最近点把向量拆成可由子空间解释的部分和正交误差。子空间条件提供中点、任意直线变化及最后的直和；闭性保留极小化列的极限，Hilbert 完备性保证该极限存在。若在 $\ell^2$ 中取非闭子空间 $c_{00}$，其正交补为零，但它不等于整个空间，所以不能删闭性。

完备性也不能一般删除。取预 Hilbert 空间 $\mathscr H=c_{00}$，$a_n=2^{-n}$，并令 $\mathscr M=\{x:\sum_nx_na_n=0\}$。Schwarz 不等式给这一定义泛函连续，故 $\mathscr M$ 在 $c_{00}$ 内闭，且 $e_1\notin\mathscr M$。若 $z\in\mathscr M^\perp$，向量 $a_me_n-a_ne_m\in\mathscr M$ 给 $a_mz_n-a_nz_m=0$。选 $m$ 超出 $z$ 的有限支集，便对每个 $n$ 得 $z_n=0$。因此 $\mathscr M^\perp=\{0\}$，直和结论失败。若只讨论某个已知完备的闭子空间，存在最近点的证明可只用该子空间完备；不能据此取消定理的普遍完备性保障。
\end{remark}
% END SUPPLEMENT: id=v14-ch05-orthogonal-decomposition-details
% SOURCE p=189; printed=174; object=ch05-inner-product-functional-introduction; region=middle
对任意 $y\in\mathscr H$，令 $f_y(x)=\langle x,y\rangle$。
% SOURCE p=189; printed=174; object=ch05-inner-product-functional-isometry; region=middle
\begin{proposition}\label{proposition:ch05-inner-product-functional-isometry}
$f_y$ 是范数为 $\|y\|$ 的有界线性泛函；映射 $y\mapsto f_y$ 是从 $\mathscr H$ 到 $\mathscr H^*$ 的共轭线性等距映射。
\end{proposition}
% SOURCE p=189; printed=174; object=ch05-riesz-introduction; region=middle
Hilbert 空间的一个基本事实是，每个有界线性泛函都这样产生。
% SUPPLEMENT: id=v14-ch05-functional-isometry-details; parent=ch05-inner-product-functional-isometry; kind=proof+motivation+hypotheses; origin=independent-derivation
\begin{proof}
\textbf{原断言补证（重构新增）。}第一变量线性给 $f_y$ 线性；Schwarz 不等式给 $|f_y(x)|\leqslant\|x\|\|y\|$，故 $\|f_y\|\leqslant\|y\|$。若 $y\ne0$，取单位向量 $x=y/\|y\|$ 就得到反向不等式；$y=0$ 时两端都为零。第二变量共轭线性给 $f_{ay+bz}=\overline a f_y+\overline b f_z$，而 $\|f_y-f_z\|=\|y-z\|$ 给等距性和单射性。
\end{proof}
\begin{remark}
\textbf{动机与条件说明（重构新增）。}每个方向 $y$ 提供一个连续的线性测量，测量强度恰等于该方向长度。此处“共轭线性”不能误写成复线性；实空间时就是线性。上述嵌入只用预 Hilbert 结构，不用完备性；下一定理的满射性才需要 Hilbert 假设。
\end{remark}
% END SUPPLEMENT: id=v14-ch05-functional-isometry-details
% SOURCE p=189; printed=174; object=ch05-riesz-hilbert-theorem; number=5.25; region=lower
\begin{theorem}\label{theorem:ch05-riesz-hilbert}
若 $f\in\mathscr H^*$，则存在唯一的 $y\in\mathscr H$，使 $f(x)=\langle x,y\rangle$ 对所有 $x\in\mathscr X$ 成立。
\end{theorem}
% SOURCE p=189; printed=174; object=ch05-riesz-hilbert-proof; region=lower
\begin{proof}
若 $\langle x,y\rangle=\langle x,y'\rangle$ 对所有 $x$ 成立，则取 $x=y-y'$ 得 $\|y-y'\|^2=0$，所以唯一性成立。若 $f=0$，取 $y=0$。否则 $\mathscr M=\ker f$ 是 $\mathscr X$ 的真闭子空间。由\cref{theorem:ch05-orthogonal-decomposition}，$\mathscr M^\perp\ne\{0\}$；取 $z\in\mathscr M^\perp$ 使 $\|z\|=1$。对每个 $x\in\mathscr H$，令 $u=f(x)z-f(z)x$，则 $u\in\mathscr M$，所以
\begin{align*}
0=\langle u,z\rangle
=f(x)\|z\|^2-f(z)\langle x,z\rangle
=f(x)-\langle x,\overline{f(z)}z\rangle.
\end{align*}
故可取 $y=\overline{f(z)}z$。
\end{proof}
% SUPPLEMENT: id=v14-ch05-riesz-hilbert-details; parent=ch05-riesz-hilbert-theorem; kind=proof-detail+source-clarification+motivation+hypotheses; origin=independent-derivation
\begin{remark}
\textbf{记号、动机与逐条件说明（重构新增）。}原定理和证明中的 $\mathscr X$ 应统一理解为 $\mathscr H$。这一定理通常称为 Hilbert 空间的 Riesz 表示定理：抽象连续线性测量能由空间内部唯一一个方向表示。$f$ 连续保证其核闭，$f\ne0$ 保证核是真子空间；正交分解因而给非零正交方向，再归一化。$f(u)=f(x)f(z)-f(z)f(x)=0$ 验证了原证的关键向量确在核中。若 $f(z)=0$，则 $z$ 同时属于 $\mathscr M$ 与 $\mathscr M^\perp$，与其单位长度矛盾。

表达式中的共轭不能省：本书内积对第二变量共轭线性，故 $\langle x,\overline{f(z)}z\rangle=f(z)\langle x,z\rangle$。上一个命题再给 $\|f\|=\|y\|$。完备性经正交分解进入证明；在 $c_{00}$ 配 $\ell^2$ 范数时，连续泛函 $f(x)=\sum_n2^{-n}x_n$ 若由 $y\in c_{00}$ 表示，代入每个 $e_n$ 就得 $y_n=2^{-n}$，矛盾。这同时说明任意连续泛函的表示不能只靠预 Hilbert 公理。反过来，任何内积表示的泛函自动连续，故一般不连续线性泛函不在结论范围内。
\end{remark}
% END SUPPLEMENT: id=v14-ch05-riesz-hilbert-details
% SOURCE p=190; printed=175; object=ch05-hilbert-reflexivity; region=top
\begin{corollary}\label{corollary:ch05-hilbert-reflexive}
Hilbert 空间在很强的意义下自反：$\mathscr H$ 不仅自然同构于 $\mathscr H^{**}$，而且与 $\mathscr H^*$ 共轭线性同构。
\end{corollary}
% SUPPLEMENT: id=v14-ch05-hilbert-reflexivity-proof; parent=ch05-hilbert-reflexivity; kind=proof+motivation+hypotheses; origin=independent-derivation
\begin{proof}
\textbf{自然双对偶识别补写（重构新增）。}记 $Vy=f_y$；已证 $V:\mathscr H\to\mathscr H^*$ 共轭线性、等距且满射。给 $F\in\mathscr H^{**}$，令 $g(y)=\overline{F(Vy)}$。两个共轭相消，故 $g$ 线性，且 $|g(y)|\leqslant\|F\|\|y\|$。由表示定理，存在 $z\in\mathscr H$ 使 $g(y)=\langle y,z\rangle$。因此
\begin{align*}
F(Vy)=\overline{g(y)}=\langle z,y\rangle=(Vy)(z).
\end{align*}
$V$ 满射使这等式覆盖所有 $f\in\mathscr H^*$，所以 $F(f)=f(z)$，即 $F=Jz$，其中 $J$ 为自然嵌入。故 $J$ 满射，空间自反。实空间中删去共轭，论证相同。
\end{proof}
\begin{remark}
\textbf{动机与条件说明（重构新增）。}自反性要求的是自然映射 $J$ 满射，不能仅说两个空间抽象同构。本证通过表示定理逐个实现双对偶泛函，明确验证了自然性。与单对偶的识别在复数情形共轭线性，与双对偶的自然识别则线性；两者应区分。完备性由表示定理继承。
\end{remark}
% END SUPPLEMENT: id=v14-ch05-hilbert-reflexivity-proof
% SOURCE p=190; printed=175; object=ch05-orthonormal-definition; region=upper
\begin{definition}[标准正交族]\label{definition:ch05-orthonormal}
$\mathscr H$ 中的族 $\{u_\alpha\}_{\alpha\in A}$ 称为\textbf{标准正交族}，若每个 $u_\alpha$ 的范数都为1，且 $\alpha\ne\beta$ 时 $u_\alpha\perp u_\beta$。
\end{definition}
% SOURCE p=190; printed=175; object=ch05-gram-schmidt-process; region=upper
\begin{definition}[Gram--Schmidt 正交化过程]\label{definition:ch05-gram-schmidt}
给定 $\mathscr H$ 中的线性无关序列 $\{x_n\}$，递归定义
\begin{align*}
u_1&=\frac{x_1}{\|x_1\|},\\
v_N&=x_N-\sum_{n=1}^{N-1}\langle x_N,u_n\rangle u_n,
\qquad u_N=\frac{v_N}{\|v_N\|}\quad(N\geqslant2).
\end{align*}
这一构造称为\textbf{Gram--Schmidt 过程}。
\end{definition}
% SOURCE p=190; printed=175; object=ch05-gram-schmidt-result; region=upper
\begin{proposition}\label{proposition:ch05-gram-schmidt}
上述过程得到标准正交序列 $\{u_n\}$，并且对每个 $N$，$\{u_1,\ldots,u_N\}$ 与 $\{x_1,\ldots,x_N\}$ 张成相同子空间。
\end{proposition}
% SOURCE p=190; printed=175; object=ch05-gram-schmidt-proof; region=upper
\begin{proof}
假设 $u_1,\ldots,u_{N-1}$ 已经构造出来。由 $\{x_n\}$ 的线性无关性，$v_N\ne0$；而对 $m<N$，
\begin{align*}
\langle v_N,u_m\rangle
=\langle x_N,u_m\rangle-\langle x_N,u_m\rangle=0.
\end{align*}
所以 $u_N=v_N/\|v_N\|$ 正是所需的第 $N$ 个向量。
\end{proof}
% SUPPLEMENT: id=v14-ch05-orthonormal-gram-schmidt-details; parent=ch05-orthonormal-definition; kind=proof-detail+motivation+hypotheses; origin=independent-derivation
\begin{proof}
\textbf{递归及张成空间补写（重构新增）。}线性无关性先给 $x_1\ne0$。假定前 $N-1$ 项结论成立，若 $v_N=0$，则 $x_N$ 属于前面 $u_n$ 的张成空间，也就是前面 $x_n$ 的张成空间，矛盾。归一化给 $\|u_N\|=1$。原正交计算利用 $\langle u_n,u_m\rangle=0$（$n\ne m$）和 $\langle u_m,u_m\rangle=1$；共轭对称性也给反方向正交。$u_N$ 是 $x_N$ 与前面 $u_n$ 的线性组合，反过来 $x_N=\|v_N\|u_N+\sum_{n<N}\langle x_N,u_n\rangle u_n$，于是两个前缀张成空间相等，完成归纳。有限个独立向量完全同样处理。
\end{proof}
\begin{remark}
\textbf{动机与逐条件说明（重构新增）。}标准化让每个方向的长度单位一致，正交让各方向互不干扰；空族也可视为标准正交族。Gram--Schmidt 每一步扣除旧方向分量，只留下新的方向。线性无关性使这个余量非零，故归一化合法；若输入相关，须先删除落在先前张成空间中的项，不能直接除以零。每一步只涉及有限和，不需要完备性；过程保留所有有限前缀张成空间，但不声称所得族已张成整个 $\mathscr H$。这里的序列条件使按自然数递归可行，不能直接用同一有限和公式处理任意不可数族。
\end{remark}
% END SUPPLEMENT: id=v14-ch05-orthonormal-gram-schmidt-details
% SOURCE p=190; printed=175; object=ch05-bessel-inequality; number=5.26; region=middle
\begin{theorem}[Bessel 不等式]\label{theorem:ch05-bessel}
若 $\{u_\alpha\}_{\alpha\in A}$ 是 $\mathscr H$ 中的标准正交族，则对每个 $x\in\mathscr H$，
\begin{align*}
\sum_{\alpha\in A}|\langle x,u_\alpha\rangle|^2\leqslant\|x\|^2.
\end{align*}
特别地，$\{\alpha:\langle x,u_\alpha\rangle\ne0\}$ 至多可数。
\end{theorem}
% SOURCE p=190; printed=175; object=ch05-bessel-proof; region=middle
\begin{proof}
只需证明任意有限 $F\subseteq A$ 上的相应不等式。由 Pythagoras 定理，
\begin{align*}
0&\leqslant\Bigl\|x-\sum_{\alpha\in F}\langle x,u_\alpha\rangle u_\alpha\Bigr\|^2\\
&=\|x\|^2-2\operatorname{Re}\Bigl\langle x,\sum_{\alpha\in F}\langle x,u_\alpha\rangle u_\alpha\Bigr\rangle
+\Bigl\|\sum_{\alpha\in F}\langle x,u_\alpha\rangle u_\alpha\Bigr\|^2\\
&=\|x\|^2-2\sum_{\alpha\in F}|\langle x,u_\alpha\rangle|^2
+\sum_{\alpha\in F}|\langle x,u_\alpha\rangle|^2
=\|x\|^2-\sum_{\alpha\in F}|\langle x,u_\alpha\rangle|^2.
\end{align*}
\end{proof}
% SUPPLEMENT: id=v14-ch05-bessel-details; parent=ch05-bessel-inequality; kind=proof-detail+motivation+hypotheses; origin=independent-derivation
\begin{proof}
\textbf{任意指标集与可数性（重构新增）。}非负和按有限部分和的上确界定义，故有限 $F$ 的共同上界就是全和的上界。对每个整数 $k\geqslant1$，令 $A_k=\{\alpha:|\langle x,u_\alpha\rangle|^2\geqslant1/k\}$。若 $A_k$ 中有任意多有限个指标，有限和界就会失败；具体地，任何有限 $F\subseteq A_k$ 都满足 $|F|/k\leqslant\|x\|^2$，所以 $A_k$ 有限。非零系数指标集等于 $\bigcup_{k\geqslant1}A_k$，故至多可数。
\end{proof}
\begin{remark}
\textbf{动机与条件说明（重构新增）。}减去有限个正交方向后的误差平方非负，所以这些方向所捕获的能量不会超过总能量。正交使交叉项消失，单位长度使系数无需再除以 $\|u_\alpha\|^2$；未归一化时直接沿用公式可失败，例如在一维取 $u=2,x=1$。指标集无需可数，结论反而推出每个向量的非零系数至多可数。完备性也未用；它将用于把平方可和的系数真正求和为向量。
\end{remark}
% END SUPPLEMENT: id=v14-ch05-bessel-details
% SOURCE p=190; printed=175; object=ch05-orthonormal-completeness; number=5.27; region=lower
\begin{theorem}\label{theorem:ch05-orthonormal-completeness}
设 $\{u_\alpha\}_{\alpha\in A}$ 为 $\mathscr H$ 中的标准正交族。以下条件等价：
\begin{enumerate}
\item\label{item:ch05-orthonormal-total} （完全性）若 $\langle x,u_\alpha\rangle=0$ 对所有 $\alpha$ 成立，则 $x=0$。
\item\label{item:ch05-parseval} （Parseval 恒等式）对每个 $x\in\mathscr H$，$\|x\|^2=\sum_{\alpha\in A}|\langle x,u_\alpha\rangle|^2$。
\item\label{item:ch05-orthogonal-expansion} 对每个 $x\in\mathscr H$，有 $x=\sum_{\alpha\in A}\langle x,u_\alpha\rangle u_\alpha$，其中只有至多可数个非零项，且无论怎样排列这些项，级数都按范数收敛。
\end{enumerate}
\end{theorem}
% SOURCE p=190,191; printed=175,176; object=ch05-orthonormal-completeness-proof; region=bottom-top
\begin{proof}
第一问蕴含第三问：将非零系数对应的指标记为 $\alpha_1,\alpha_2,\ldots$。Bessel 不等式说明 $\sum_j|\langle x,u_{\alpha_j}\rangle|^2$ 收敛；由 Pythagoras 定理，
\begin{align*}
\Bigl\|\sum_{j=n}^m\langle x,u_{\alpha_j}\rangle u_{\alpha_j}\Bigr\|^2
=\sum_{j=n}^m|\langle x,u_{\alpha_j}\rangle|^2\longrightarrow0
\qquad(n,m\to\infty).
\end{align*}
因此 $\sum_j\langle x,u_{\alpha_j}\rangle u_{\alpha_j}$ 由 $\mathscr H$ 的完备性收敛。若令 $y=x-\sum_j\langle x,u_{\alpha_j}\rangle u_{\alpha_j}$，显然对所有 $\alpha$ 都有 $\langle y,u_\alpha\rangle=0$，于是由第一问，$y=0$。

第三问蕴含第二问：沿用上述记号，像 Bessel 不等式的证明一样，有
\begin{align*}
\|x\|^2-\sum_{j=1}^n|\langle x,u_{\alpha_j}\rangle|^2
=\Bigl\|x-\sum_{j=1}^n\langle x,u_{\alpha_j}\rangle u_{\alpha_j}\Bigr\|^2\longrightarrow0.
\end{align*}
最后，第二问蕴含第一问是显然的。
\end{proof}
% SUPPLEMENT: id=v14-ch05-orthogonal-expansion-details; parent=ch05-orthonormal-completeness; kind=proof-detail+motivation+hypotheses; origin=independent-derivation
\begin{proof}
\textbf{排列、极限与边界情形（重构新增）。}若非零系数有限，直接取有限和 $s$，由正交性得 $\langle x-s,u_\alpha\rangle=0$，于是完全性给 $x=s$；若没有非零系数，则直接给 $x=0$。无限情形中，对固定 $\alpha$，部分和与 $u_\alpha$ 的内积或者始终为0，或者在该指标出现后恒为 $\langle x,u_\alpha\rangle$。内积连续，所以原证明的余量与每个 $u_\alpha$ 都正交。

非负系数平方和不依赖排列：给 $\varepsilon>0$，按有限部分和上确界可选有限指标集 $F$，使 $F$ 外的平方和小于 $\varepsilon^2$。任何排列在包含 $F$ 后的任意有限尾部，其对应向量范数小于 $\varepsilon$。因此每种排列的部分和都 Cauchy；完备性给其极限，同一个完全性论证又使每种极限都等于 $x$。这证明了原文的“无论怎样排列”。第三问到第二问用有限误差恒等式后取极限；第二问到第一问则因全部系数为零给 $\|x\|^2=0$，再由正定性得 $x=0$。
\end{proof}
\begin{remark}
\textbf{动机与逐条件说明（重构新增）。}完全性表示没有任何非零方向被这族测量遗漏；Parseval 表示没有能量遗漏；展开式表示向量本身能由测量重建。三种说法对应同一事实。标准正交性负责系数公式和精确误差平方，空间完备性在第一问推出级数存在时使用。此处族的“完全性”与空间的“完备性”是两个概念。指标集可不可数，但每个具体展开的非零支集至多可数；这不要求整个空间可分。零空间的空标准正交族满足三个条件。所证为正交级数的无条件收敛，不等于 $\sum_\alpha|\langle x,u_\alpha\rangle|<\infty$，例如 $\ell^2$ 中系数 $1/n$ 就不绝对可和。
\end{remark}
% END SUPPLEMENT: id=v14-ch05-orthogonal-expansion-details
% SOURCE p=191; printed=176; object=ch05-orthonormal-basis-definition; region=upper
\begin{definition}[标准正交基]\label{definition:ch05-orthonormal-basis}
满足\cref{theorem:ch05-orthonormal-completeness}中三个等价条件的标准正交族，称为 $\mathscr H$ 的\textbf{标准正交基}。
\end{definition}
% SOURCE p=191; printed=176; object=ch05-ell2-standard-basis-example; region=upper
\begin{example}\label{example:ch05-ell2-standard-basis}
设 $\mathscr H=\ell^2(A)$。对每个 $\alpha\in A$，定义 $e_\alpha\in\ell^2(A)$：当 $\beta=\alpha$ 时 $e_\alpha(\beta)=1$，否则 $e_\alpha(\beta)=0$。族 $\{e_\alpha\}_{\alpha\in A}$ 显然标准正交，而且对任意 $f\in\ell^2(A)$ 有 $\langle f,e_\alpha\rangle=f(\alpha)$，所以它是一个标准正交基。
\end{example}
% SUPPLEMENT: id=v14-ch05-orthonormal-basis-details; parent=ch05-orthonormal-basis-definition; kind=proof-detail+motivation+hypotheses; origin=independent-derivation
\begin{remark}
\textbf{动机与定义边界（重构新增）。}这里的“基”允许无穷范数收敛展开，区别于只允许有限线性组合的代数基。标准正交族未必是基，例如 $\ell^2$ 中删去 $e_1$ 后就遗漏第一坐标。原例中每个 $e_\alpha$ 的平方和为1，不同坐标的乘积处处为0；因此标准正交。若 $f$ 与所有 $e_\alpha$ 正交，逐坐标公式给 $f(\alpha)=0$ 对所有 $\alpha$，故 $f=0$，严格验证了完全性。任意集合 $A$ 均可使用，不需要预先编号为一个序列。
\end{remark}
% END SUPPLEMENT: id=v14-ch05-orthonormal-basis-details
% SOURCE p=191; printed=176; object=ch05-orthonormal-basis-existence; number=5.28; region=middle
\begin{proposition}\label{proposition:ch05-orthonormal-basis-existence}
每个 Hilbert 空间都有标准正交基。
\end{proposition}
% SOURCE p=191; printed=176; object=ch05-orthonormal-basis-existence-proof; region=middle
\begin{proof}
通常地应用 Zorn 引理可知，所有标准正交集按包含关系组成的集合有极大元；极大性等价于\subcref{theorem:ch05-orthonormal-completeness}{item:ch05-orthonormal-total}。
\end{proof}
% SUPPLEMENT: id=v14-ch05-orthonormal-basis-zorn-details; parent=ch05-orthonormal-basis-existence; kind=proof-detail+motivation+hypotheses; origin=independent-derivation
\begin{proof}
\textbf{Zorn 论证完整补写（重构新增）。}令 $\mathscr P$ 为 $\mathscr H$ 的所有标准正交子集，包含空集，按包含排序。它是幂集的子集，故是集合，且非空。对任意非空链，取其并集；并中任意两个向量同属于链中较大的某个成员，故各为单位向量且不同者正交。因此并集仍在 $\mathscr P$，并为链的上界。空链以空集为上界。由\hyperref[lemma:ch00-zorn]{Zorn 引理}得极大成员 $U$。若非零 $x$ 与 $U$ 中每个向量正交，则 $x/\|x\|$ 不属于 $U$，却可添入而仍标准正交，矛盾。故 $U$ 完全，是标准正交基。反之，完全的标准正交族不能再添单位正交向量，故极大。若 $\mathscr H=\{0\}$，所得基就是空集。
\end{proof}
\begin{remark}
\textbf{动机与条件说明（重构新增）。}极大化把“不断添加新正交方向”推广到不可数维情形；每条链的并仍满足两两条件，使 Zorn 引理适用。这里使用选择原理；它给存在性，不给具体计算算法。Zorn 步骤本身只用内积，Hilbert 完备性通过先前的等价定理保证极大族的无穷展开性质。
\end{remark}
% END SUPPLEMENT: id=v14-ch05-orthonormal-basis-zorn-details
% SOURCE p=191; printed=176; object=ch05-separable-basis; number=5.29; region=middle
\begin{proposition}\label{proposition:ch05-separable-orthonormal-basis}
Hilbert 空间 $\mathscr H$ 可分，当且仅当它有可数标准正交基；此时它的每个标准正交基都可数。
\end{proposition}
% SOURCE p=191; printed=176; object=ch05-separable-basis-proof; region=lower
\begin{proof}
若 $\{x_n\}$ 是 $\mathscr H$ 的可数稠密集，递归地删去属于 $x_1,\ldots,x_{n-1}$ 线性张成空间的 $x_n$，就得到线性无关序列 $\{y_n\}$，其张成空间在 $\mathscr H$ 中稠密。对它应用 Gram--Schmidt 过程，得到标准正交序列 $\{u_n\}$，其张成空间也稠密，因而是基。反过来，若 $\{u_n\}$ 为可数标准正交基，取复数域中某个可数稠密子集，以其中元素为系数的有限线性组合构成 $\mathscr H$ 的一个可数稠密集。

此外，若 $\{v_\alpha\}_{\alpha\in A}$ 为另一个标准正交基，则对每个 $n$，$A_n=\{\alpha\in A:\langle u_n,v_\alpha\rangle\ne0\}$ 可数。由 $\{u_n\}$ 的完全性，$A=\bigcup_{n=1}^{\infty}A_n$，所以 $A$ 可数。
\end{proof}
% SUPPLEMENT: id=v14-ch05-separable-basis-details; parent=ch05-separable-basis; kind=proof-detail+motivation+hypotheses; origin=independent-derivation
\begin{proof}
\textbf{稠密性与有限维边界（重构新增）。}递归只保留不在此前已保留向量张成空间中的项。每个被删项属于此前保留项的有限张成空间，因此全部原稠密集仍在保留项的张成空间内。保留项可能有限或为空；有限时使用有限 Gram--Schmidt，空时原空间只能为零。若 $x$ 与所得标准正交族都正交，则与其张成空间正交；连续性把这延伸到其闭包即整个空间，取内积对象为 $x$ 得 $x=0$。

反向可取系数集 $\mathbb Q+\mathrm i\mathbb Q$（实空间取 $\mathbb Q$），再加零组合。有限指标选择和有限个可数系数的选择总共可数。给 $x$ 和 $\varepsilon>0$，先用正交展开选有限部分和使误差小于 $\varepsilon/2$；若该和有 $N>0$ 项，逐项把系数逼近到误差小于 $\varepsilon/(2\sqrt N)$，Pythagoras 定理使系数替换误差小于 $\varepsilon/2$。空和无须替换，因此该可数集合稠密。最后，Bessel 不等式给每个 $A_n$ 可数；任意 $v_\alpha$ 范数为1，不可能与完整的 $\{u_n\}$ 全正交，故确有 $A=\bigcup_nA_n$。有限及空的首个基同理。
\end{proof}
\begin{remark}
\textbf{动机与条件说明（重构新增）。}可分性把全空间的逼近归结为可数数据，正交化将这些数据整理成独立方向。“可数”包括有限和空的情形，不能为了使用序列而错误推出零空间有单位向量。反向不只用有理实系数逼近复系数，还要允许有理虚部。每个基都可数的结论由 Bessel 与完全性得到，不是任意代数基的结论。
\end{remark}
% END SUPPLEMENT: id=v14-ch05-separable-basis-details
% SOURCE p=191; printed=176; object=ch05-separable-examples-reference; region=lower
实际中出现的大多数 Hilbert 空间都是可分的。我们将在\cref{exercise:ch05-l2-disjoint-decomposition,exercise:ch05-product-orthonormal-basis,exercise:ch05-lebesgue-l2-separable}中讨论一些例子。
% SOURCE p=191; printed=176; object=ch05-unitary-definition; region=lower
\begin{definition}[酉映射]\label{definition:ch05-unitary-map}
若 $\mathscr H_1,\mathscr H_2$ 为内积分别记作 $\langle\cdot,\cdot\rangle_1,\langle\cdot,\cdot\rangle_2$ 的 Hilbert 空间，从 $\mathscr H_1$ 到 $\mathscr H_2$ 的\textbf{酉映射}是一个可逆线性映射 $U:\mathscr H_1\to\mathscr H_2$，它保持内积：
\begin{align*}
\langle Ux,Uy\rangle_2=\langle x,y\rangle_1\qquad(x,y\in\mathscr H_1).
\end{align*}
\end{definition}
% SOURCE p=191; printed=176; object=ch05-unitary-isometry-discussion; region=bottom
取 $y=x$，可见每个酉映射都是等距映射：$\|Ux\|_2=\|x\|_1$。反过来，每个满的等距映射都是酉映射，见\cref{exercise:ch05-polarization-unitary}。酉映射是 Hilbert 空间范畴中真正的“同构”：它不仅保持线性结构和拓扑，还保持范数和内积。从这种抽象结构的角度看，每个 Hilbert 空间都形同一个 $\ell^2$ 空间。
% SUPPLEMENT: id=v14-ch05-unitary-definition-details; parent=ch05-unitary-definition; kind=source-clarification+motivation+hypotheses; origin=independent-derivation
\begin{remark}
\textbf{线性条件与动机（重构新增）。}原段反向断言须在\subcref{exercise:ch05-polarization-unitary}{item:ch05-unitary-isometric-surjective}规定的线性映射范围内理解。一般满的度量等距映射未必线性，例如 $z\mapsto z+1$；即使固定0，$\mathbb C$ 上 $z\mapsto\overline z$ 也是满等距映射，却不复线性且不保持本书的复内积。在线性前提下，极化恒等式从范数恢复内积，完整证明见所引习题。

保持内积已蕴含单射，因为 $Ux=0$ 给 $\|x\|=0$；可逆性额外要求满射，排除只嵌入目标真子空间的情形，例如 $\ell^2$ 上的右移。它还自动给逆映射等距连续，所以无须在定义中再加有界性。Hilbert 空间是定义的使用背景；线性满等距与保内积之间的这个代数判别本身不需要完备性。
\end{remark}
% END SUPPLEMENT: id=v14-ch05-unitary-definition-details
% SOURCE p=192; printed=177; object=ch05-hilbert-ell2-unitary; number=5.30; region=top
\begin{proposition}\label{proposition:ch05-hilbert-ell2-unitary}
设 $\{e_\alpha\}_{\alpha\in A}$ 为 $\mathscr X$ 的一个标准正交基。由 $\widehat x(\alpha)=\langle x,u_\alpha\rangle$ 定义的对应 $x\mapsto\widehat x$ 是从 $\mathscr H$ 到 $\ell^2(A)$ 的酉映射。
\end{proposition}
% SOURCE p=192; printed=177; object=ch05-hilbert-ell2-unitary-proof; region=upper
\begin{proof}
映射 $x\mapsto\widehat x$ 显然线性；由 Parseval 恒等式 $\|x\|^2=\sum|\widehat x(\alpha)|^2$，它是从 $\mathscr H$ 到 $\ell^2(A)$ 的等距映射。若 $f\in\ell^2(A)$，则 $\sum|f(\alpha)|^2<\infty$，所以 Pythagoras 定理表明级数 $\sum f(\alpha)u_\alpha$ 的部分和为 Cauchy 列，其中只有至多可数项非零。因此 $x=\sum f(\alpha)u_\alpha$ 存在于 $\mathscr H$，且 $\widehat x=f$。由\subcref{exercise:ch05-polarization-unitary}{item:ch05-unitary-isometric-surjective}，$x\mapsto\widehat x$ 为酉映射。
\end{proof}
% SUPPLEMENT: id=v14-ch05-hilbert-ell2-details; parent=ch05-hilbert-ell2-unitary; kind=proof-detail+source-clarification+motivation+hypotheses; origin=independent-derivation
\begin{remark}
\textbf{原记号统一（重构新增）。}原命题先写基为 $\{e_\alpha\}$、空间为 $\mathscr X$，公式及证明却用 $\{u_\alpha\}$、$\mathscr H$。这里应始终指同一个 Hilbert 空间 $\mathscr H$ 及同一组标准正交基 $\{u_\alpha\}$；并非引入另一空间或另一基。帽号在此表示正交坐标变换，应保留其语义，不改成弧帽。
\end{remark}
\begin{proof}
\textbf{满射和坐标恢复补写（重构新增）。}Bessel 不等式先保证每个 $\widehat x$ 确在 $\ell^2(A)$；第一变量线性给该映射复线性，Parseval 再给等距。对 $f\in\ell^2(A)$，按非负平方和的有限性，其非零支集至多可数；若支集有限或空，只取相应有限和。否则枚举支集，有限尾和的平方范数恰为系数平方尾和，趋零，故由完备性存在极限 $x$。对任意 $\beta$，部分和与 $u_\beta$ 的内积在含该项后等于 $f(\beta)$；不在支集时始终为0。由内积连续性得 $\langle x,u_\beta\rangle=f(\beta)$，故满射。极化习题只用内积公理，不依赖本命题，引用它没有循环。
\end{proof}
\begin{remark}
\textbf{动机与条件说明（重构新增）。}选择标准正交基后，每个向量等价于一组平方可和坐标，且长度和内积都保持。正交性负责坐标的独立性，标准化负责单位权重，完全性负责等距和无遗漏，空间完备性负责给定平方可和系数的合成。无需可分性；不可数 $A$ 允许不可分空间，只是每个向量仍具有可数支集。映射依赖所选基，不能称为不经选择的唯一坐标识别。
\end{remark}
% END SUPPLEMENT: id=v14-ch05-hilbert-ell2-details
\subfile{section-05-exercises.tex}
\end{document}
